% SaT-CPS 2026 submission. ACM proceedings, sigconf, 10 pages.
% Use only ACM-approved packages; no manual margin or font overrides.
\documentclass[sigconf]{acmart}

\title{Secure CPS Tool Invocation via a Typed-Plan Firewall}

% Author names and affiliations per workshop CFP; replace placeholders
\author{Author One}
\affiliation{%
  \institution{Institution One}
  \city{City}
  \country{Country}
}
\email{author1@example.org}

\author{Author Two}
\affiliation{%
  \institution{Institution Two}
  \city{City}
  \country{Country}
}
\email{author2@example.org}

% CCS concepts (ACM style)
\ccsdesc[500]{Security and privacy~Security in hardware}
\ccsdesc[300]{Computer systems organization~Embedded and cyber-physical systems}
\ccsdesc[300]{Software and its engineering~Software safety}

\keywords{cyber-physical systems, security, runtime enforcement, tool invocation, typed plans, access control, auditability}

\begin{document}

\begin{abstract}
Tool invocation in cyber-physical control planes is a trust boundary: planner output can request unsafe tools or arguments (e.g., path traversal), and execution without validation undermines safety and auditability. We describe a runtime enforcement architecture for CPS tool invocation: typed plans, tool allow-listing, argument validation (safe\_args + ponr\_gate), and deterministic capture of proposed invocations. The validator blocks all released synthetic unsafe cases and admits the released safe cases (15 red-team, 6 confusable deputy, 4 jailbreak-style). Real-LLM evaluation on OpenAI in the camera-ready snapshot (\texttt{llm\_eval\_camera\_ready\_20260424}, full-suite mode, 3 runs/case, 25 cases/model): \texttt{gpt-4.1-mini} and \texttt{gpt-4.1} each achieve 75/75 (100.0\%, 95\% Wilson CI [95.1, 100.0\%]). An optional four-model matrix via Prime Inference is documented in the venue runbook and uses a separate denominator. The security layer adds bounded measured overhead; denial traces are recorded for replay and audit. Tool-level baseline: gated and weak both deny the injected disallowed tool; ungated allows. Argument-level baseline (safe\_args ablation): gated denies, weak and ungated allow. Benign-suite results (optional) quantify false positives on safe steps. We claim containment and auditability; we do not claim elimination of prompt injection or full coverage of adversarial inputs.
\end{abstract}

\maketitle

\section{Introduction}

\paragraph{CPS tool invocation as a security problem.}
In CPS, the control plane often invokes tools (query status, submit results, or legacy interfaces). When a planner---whether an LLM or a heuristic---produces tool calls, those calls cross a trust boundary: unchecked execution of disallowed tools or malicious arguments (e.g., path traversal) can compromise safety and integrity. Runtime enforcement at this boundary is therefore a core concern for secure CPS architectures.

\paragraph{Why generic LLM framing is insufficient.}
Generic LLM safety benchmarks focus on prompt injection, jailbreaks, or output filtering in isolation. For CPS we need: (a) a clear trust boundary (planner output is untrusted until validated); (b) integration with trace and audit (deterministic capture, denial logs); (c) evaluation that includes both synthetic validator correctness and real-LLM stress tests; (d) explicit baselines that show the value of each validation layer (allow-list vs.\ safe\_args).

\paragraph{The typed-plan firewall.}
The typed-plan firewall sits at the actuation boundary: planner output is untrusted until the firewall admits it. Plans are typed (tool, args, validators); the firewall runs allow-list and argument checks (e.g., normalized path traversal, dangerous patterns, recursive denylist-key checks). In the published mock execution harness, denied steps are not executed and are captured for replay and audit. This is a runtime enforcement architecture for CPS tool invocation, not merely an LLM benchmark.

\paragraph{Evidence and scope.}
We report four evaluation blocks: (A) synthetic validator evidence; (B) real-LLM stress test with explicit failure-case analysis; (C) adapter latency and denial traces; (D) baseline comparison (tool-level and argument-level). We do not claim elimination of prompt injection, full jailbreak coverage, or universal robustness.

\section{Threat Model}

\textbf{Assets.} CPS tool interfaces; operational state accessible via those tools; traces and denial logs; policy configuration (allow-list, validator stack).

\textbf{Adversarial leverage.} Prompt injection via planner input; tool-returned content that induces privilege-seeking behavior; malformed or privilege-seeking arguments (e.g., path traversal); confused-deputy dynamics; jailbreak-style instruction phrases in arguments.

\textbf{Trust assumptions.} The planner is untrusted at the actuation boundary. The validator and capture pipeline are trusted. The execution environment obeys firewall decisions (denied steps are not executed).

\textbf{Security properties sought.} No unsafe step (disallowed tool or unsafe arguments in the released suite) reaches execution in the reported setting; all denied actions are captured for replay and audit; runtime overhead remains low enough for the target CPS context.

\textbf{Non-goals.} Proving semantic correctness of whole plans; complete jailbreak coverage; universal robustness; elimination of prompt injection.

\begin{table}[t]
\caption{Threat model summary.}
\label{tab:threat}
\centering
\small
\begin{tabular}{@{}lllll@{}}
\hline
Threat & Entry & Property & Mitigation & Out of scope \\
\hline
Disallowed tool & Planner & No unsafe step & allow\_list & Semantic correctness \\
Path traversal & Planner & No unsafe step & safe\_args & Full sanitization \\
Jailbreak-style args & Planner & No unsafe step & safe\_args & General injection \\
Confused deputy & Planner & No unsafe step & safe\_args & Cross-component trust \\
Unaudited execution & Adapter & Replay/audit & Capture & Full forensics \\
\hline
\end{tabular}
\end{table}

\section{System Overview}

Typed plans use a schema with steps (seq, tool, args, validators). The validator stack includes \texttt{allow\_list} (tool allow-listing) and \texttt{safe\_args} (path traversal, dangerous patterns, jailbreak-style phrases). \texttt{validate\_plan\_step} runs the configured validators; \texttt{policy\_check\_step} provides allow-list-only (weak) mode. Denied steps are recorded in trace metadata; \texttt{capture\_tool\_call} supports replay. The LLMPlanningAdapter integrates with MAESTRO: it runs scenarios, injects plan and validation outcome into trace, and reports denials and tasks completed.
% Figure 1 (decision path): python scripts/export_p6_firewall_flow.py -> docs/figures/p6_firewall_flow.mmd; render to PDF/PNG and include as fig1.pdf or fig1.png
% \begin{figure}[t]
% \centering
% \includegraphics[width=\columnwidth]{fig1}
% \caption{Decision path: planner output to allow/deny (typed step -> allow-list -> safe\_args -> capture).}
% \label{fig:decisionpath}
% \end{figure}

\section{Evaluation}

\subsection{Block A: Synthetic Validator Evidence}

Synthetic evidence isolates validator correctness from model-generation variability. The full synthetic table (15 red-team, 6 confusable deputy, 4 jailbreak-style) is the primary validator evidence: all unsafe cases blocked, all safe cases admitted.

\subsection{Block B: Real-LLM Stress Test}

The validator correctly blocks the canonical unsafe step; the real-LLM experiment verifies that, when two OpenAI models produce the targeted typed unsafe forms under the controlled full-suite prompt, the validator yields the expected allow/deny outcomes. \textbf{OpenAI (camera-ready snapshot):} full-suite mode, 3 runs/case, 25 cases/model (75 trials/model). \texttt{gpt-4.1-mini} and \texttt{gpt-4.1} each score 75/75 (100.0\%, 95\% Wilson CI [95.1, 100.0\%]). \textbf{Optional Prime matrix:} separate denominator and output directory; see venue \texttt{EXPERIMENTS\_RUNBOOK.md} and do not merge rows without relabeling.

\paragraph{Failure case as first-class result.} Any future argument-level misses or parse failures must be reported explicitly with denominator and suite mode. Repeated-run real-LLM evaluation is necessary to assess containment; argument-level cases are exactly where real systems need tighter controls.

\subsection{Block C: Adapter Latency and Denial Traces}

Three scenarios, 20 seeds each; mean p95 latency, stdev, 95\% CI; total denied unsafe steps; per-scenario denial counts; tasks completed. This security layer adds bounded measured overhead while preserving task completion in the reported denial-injection setting.

\paragraph{Case study: denial trace.} In a representative run, the planner proposed step \texttt{\{ tool: execute\_system, args: \{ cmd: ``rm -rf'' \} \}}. Allow-list check: tool not in allow\_list $\rightarrow$ Deny. The step was captured for audit; no execution occurred.

\subsection{Block D: Baseline Comparison}

\textbf{Tool-level benchmark.} The injected unsafe step is \texttt{execute\_system} (disallowed tool). Gated and weak both deny (rejection at allow-list layer); ungated allows. This baseline shows the value of having a gate at all; it does not isolate the marginal contribution of safe\_args and does not replace required nominal steps.

\textbf{Argument-level benchmark (safe\_args ablation).} With an allow-listed tool and unsafe arguments (path traversal), gated denies (safe\_args blocks), weak allows (allow-list only), ungated allows. This demonstrates the incremental value of argument-level validation. Additional task-critical replacement/competition evidence is captured in \texttt{task\_critical\_injection.json}.

\textbf{Benign utility (optional).} All-safe plans across gated/weak/ungated quantify false-positive denials; see \texttt{baseline\_benign.json} and \texttt{p6\_benign\_suite.json}.

\section{Discussion and Limitations}

Containment, not elimination: the firewall blocks the released unsafe cases; we do not claim elimination of prompt injection or full coverage. \texttt{safe\_args} is deterministic deny-pattern logic (recursive denylist-key checks, normalized path-traversal checks, and substring checks), not a full per-tool semantic sanitizer. The mock execution harness validates deny-vs-execute behavior for this artifact; production external-tool execution semantics remain future work. Real-LLM pass rate is model- and prompt-dependent; synthetic table remains primary validator evidence. Residual risk from novel attack vectors remains.

\section{Conclusion}

We presented a runtime enforcement architecture for CPS tool invocation: typed-plan firewall with allow-listing and argument validation, deterministic capture for audit, and a released evaluation suite (synthetic, real-LLM, adapter latency, tool-level and argument-level baselines). The firewall blocks the released unsafe cases and admits the released safe cases; denial-injection experiments show bounded overhead and the incremental value of argument-level validation. We report the real-LLM path-traversal failure case as a first-class result. Claims are limited to containment and auditability under the reported evaluation.

\bibliographystyle{ACM-Reference-Format}
\bibliography{references}

\end{document}

